\begin{frame}{КФ как средний вес пары}
    \small
    $A(\mathbf q)$ --- распределение пар в числителе,
    $B(\mathbf q)$ --- опорное распределение; $C=\mathcal{N}A/B$.
    $\mathcal{N}$ --- нормировка, $n$ --- число пар в ячейке.
    Здесь $\mathcal{N}=1$, обе гистограммы заполняются по одним $n$ парам:
    \[
        A = \sum_{i=1}^{n} w_i, \qquad
        B = \sum_{i=1}^{n} 1 = n, \qquad Q_A = \sum_{i=1}^{n} w_i^2.
    \]
    Поэтому КФ в одной ячейке представляет собой средний вес:
    \[
        C = \frac{A}{B} = \frac{A}{n} = \overline{w}, \qquad n>0.
    \]
    \begin{itemize}
        \item $Q_A$ хранится в \texttt{Sumw2} взвешенной гистограммы.
        \item При случайном $n$ общие пары могут связывать суммы $A$ и $B$.
    \end{itemize}
\end{frame}

\begin{frame}{Ошибка среднего при фиксированном числе пар}
    \small
    Пусть $n$ фиксировано, а веса независимы и одинаково распределены
    с дисперсией $\sigma_w^2$. Тогда
    \[
        \operatorname{Var}(C\mid n) = \frac{\sigma_w^2}{n}.
    \]
    Несмещённая выборочная оценка при $n>1$:
    \[
        s_w^2 = \frac{Q_A-A^2/n}{n-1}, \qquad
        \widehat{\operatorname{Var}}(C\mid n)
        = \frac{s_w^2}{n}
        = \frac{Q_A-A^2/n}{n(n-1)}.
    \]
    \begin{itemize}
        \item $\sigma_w^2$ --- дисперсия распределения весов;
              $s_w^2$ --- её оценка по наблюдаемым парам.
        \item Поправка $n-1$ учитывает оценивание среднего по той же выборке.
    \end{itemize}
\end{frame}

\begin{frame}{Частный случай: бинарный отбор}
    \small
    Для индикатора отбора $I_i\in\{0,1\}$ выполнено $I_i^2=I_i$.
    Поэтому $Q_A=A$, а $\widehat p=A/B=A/n$ оценивает вероятность отбора $p$.
    \[
        \operatorname{Var}(\widehat p\mid n) = \frac{p(1-p)}{n}.
    \]
    Две оценки этой дисперсии:
    \[
        \begin{aligned}
            \widehat V_{\mathrm{ROOT\ B}}
                &= \frac{\widehat p(1-\widehat p)}{n},
                &&\text{подстановка }p\to\widehat p;\\
            \widehat V_{\texttt{cf\_maker}}
                &= \frac{\widehat p(1-\widehat p)}{n-1},
                &&\text{несмещённая оценка},\quad n>1.
        \end{aligned}
    \]
    Одного ограничения $0\leq w_i\leq 1$ недостаточно:
    для небинарных весов равенство $Q_A=A$ обычно не выполняется.
\end{frame}

\begin{frame}{Общие суммы: пуассоновская модель}
    \small
    Здесь число вкладов случайно: $n\sim\operatorname{Pois}(\mu)$.
    Вклады $(a_i,b_i)$ независимы, одинаково распределены и не зависят от $n$.
    \[
        A=\sum_i a_i, \qquad B=\sum_i b_i, \qquad C=\frac{A}{B}.
    \]
    Суммы для оценки дисперсий и ковариации:
    \[
        Q_A=\sum_i a_i^2, \qquad
        Q_B=\sum_i b_i^2, \qquad Q_{AB}=\sum_i a_i b_i.
    \]
    $\widehat V_A=Q_A$, $\widehat V_B=Q_B$,
    $\widehat{\operatorname{Cov}}(A,B)=Q_{AB}$.
    Разложение отношения до первого порядка даёт
    \[
        \widehat{\operatorname{Var}}(C)\simeq
        \frac{Q_A-2C Q_{AB}+C^2 Q_B}{B^2},
        \qquad B\ne 0.
    \]
    Это дельта-оценка, а не точная дисперсия отношения для малой выборки.
\end{frame}

\begin{frame}{Ковариация взвешенной суммы и числа пар}
    \small
    Для КФ в пуассоновской модели $a_i=w_i$, $b_i=1$. Поэтому
    \[
        \widehat V_A=Q_A, \qquad \widehat V_B=n, \qquad
        \widehat{\operatorname{Cov}}(A,B)=\sum_i w_i\cdot 1=A.
    \]
    Подставляя $C=A/B$ и $B=n$ в дельта-оценку, получаем
    \[
        \widehat V_{\Delta}=
        \frac{Q_A-2CA+C^2B}{B^2}
        =\frac{Q_A-A^2/n}{n^2}.
    \]
    Оценка \texttt{cf\_maker} для фиксированного $n$ отличается поправкой:
    \[
        \widehat{\operatorname{Var}}(C\mid n)
        =\frac{n}{n-1}\,\widehat V_{\Delta}, \qquad n>1.
    \]
    При большом $n$ оценки асимптотически совпадают.
    Их статистические допущения различаются.
\end{frame}

\begin{frame}{КФ и взвешенная эффективность: разные вклады}
    \small
    Для общих пар различается связь числителя и знаменателя:
    \[
        \begin{array}{l|c|c|c}
            & a_i & b_i & \widehat{\operatorname{Cov}}(A,B)\\ \hline
            \text{КФ} & w_i & 1 & A\\
            \text{эффективность} & u_i I_i & u_i & Q_{\mathrm{pass}}
        \end{array}
    \]
    Для эффективности $\varepsilon=A/B$, $I_i\in\{0,1\}$, $u_i\geq 0$ и
    \[
        Q_{\mathrm{pass}}=\sum_i u_i^2 I_i=\widehat V_A,
        \qquad Q_{\mathrm{all}}=\sum_i u_i^2.
    \]
    Дельта-оценка эффективности имеет вид формулы ROOT \texttt{B}:
    \[
        \widehat{\operatorname{Var}}(\varepsilon)\simeq
        \frac{(1-2\varepsilon)Q_{\mathrm{pass}}+\varepsilon^2 Q_{\mathrm{all}}}{B^2}.
    \]
    Произвольные веса $w_i$ в числителе не обеспечивают эту ковариацию.
    \par\medskip
    {\footnotesize Источник:
        \href{https://root.cern.ch/doc/master/TH1_8cxx_source.html}{ROOT: TH1::Divide, опция B}.}
\end{frame}

\begin{frame}{Контрольный пример: произвольные веса}
    \small
    По $n=100$ пар, $B=n$. Проверка: ROOT 6.40.02.
    \begin{center}
        \renewcommand{\arraystretch}{1.35}
        \begin{tabular}{lccc}
            \hline
            Веса & $C$ & $\sigma_C$, \texttt{cf\_maker} & $\sigma_C$, ROOT \texttt{B}\\ \hline
            Все $w_i=1$ & 1 & 0 & 0\\
            50 весов 0, 50 весов 2 & 1 & 0.1005 & 0\\
            Все $w_i=2$ & 2 & 0 & 0.2828\\ \hline
        \end{tabular}
        \par\smallskip
        Контрольный пример, не данные столкновений.
    \end{center}
    \begin{itemize}
        \item Во второй строке $A=B$, но веса изменяются, и ошибка среднего ненулевая.
        \item ROOT \texttt{B} задаёт нулевую ошибку при $A=B$.
              Это не проверка применимости модели эффективности.
    \end{itemize}
\end{frame}

\begin{frame}{Корреляции пар и применимость ошибок}
    \small
    Число пар $n$ фиксировано, но веса могут быть зависимы:
    \[
        \operatorname{Var}(C\mid n)=\frac{1}{n^2}
        \left[\sum_i\operatorname{Var}(w_i\mid n)
        +2\sum_{i<j}\operatorname{Cov}(w_i,w_j\mid n)\right].
    \]
    \begin{itemize}
        \item Общие события или треки могут связывать веса разных пар.
        \item Bootstrap или jackknife по событиям сохраняют
              совместные вклады пар внутри события.
        \item При $n\leq 1$ выборочная дисперсия недоступна;
              нулевая записанная ошибка не означает точное измерение.
    \end{itemize}
    Опция ROOT \texttt{B} не устраняет событийные корреляции.
    Для интерпретации $\chi^2$ нужно проверить допущения об ошибках.
\end{frame}
